%% Tables and numbers in this section are generated by
%%   python replication/bank_replication.py v113i06.dta W.csv
%% from the replication materials of Kripfganz and Sarafidis (2025).
\providecommand{\BankMaxRelCoef}{}\renewcommand{\BankMaxRelCoef}{0.000}
\providecommand{\BankMaxRelSE}{}\renewcommand{\BankMaxRelSE}{0.00}
\providecommand{\BankMaxMGdiffSE}{}\renewcommand{\BankMaxMGdiffSE}{0.00}
\providecommand{\BankMedMGdiffSE}{}\renewcommand{\BankMedMGdiffSE}{0.00}
\providecommand{\BankJ}{}\renewcommand{\BankJ}{18.825}
\providecommand{\BankpJ}{}\renewcommand{\BankpJ}{0.468}
\providecommand{\BankNflag}{}\renewcommand{\BankNflag}{95}
\providecommand{\BankNunits}{}\renewcommand{\BankNunits}{350}
\providecommand{\BankNflagSize}{}\renewcommand{\BankNflagSize}{24}
\providecommand{\BankNflagBuffer}{}\renewcommand{\BankNflagBuffer}{15}
\providecommand{\BankNill}{}\renewcommand{\BankNill}{5}
\providecommand{\BankNillFlag}{}\renewcommand{\BankNillFlag}{1}
\providecommand{\BankMGpsi}{}\renewcommand{\BankMGpsi}{0.032}
\providecommand{\BankMGpsiNoDD}{}\renewcommand{\BankMGpsiNoDD}{0.151}
\providecommand{\BankMGsize}{}\renewcommand{\BankMGsize}{2.004}
\providecommand{\BankMGbuffer}{}\renewcommand{\BankMGbuffer}{$-$0.376}
\providecommand{\BankTwoSize}{}\renewcommand{\BankTwoSize}{0.223}
\providecommand{\BankTwoBuffer}{}\renewcommand{\BankTwoBuffer}{$-$0.055}
\providecommand{\BankTwoPsi}{}\renewcommand{\BankTwoPsi}{0.394}
\providecommand{\BankMedSize}{}\renewcommand{\BankMedSize}{0.807}
\providecommand{\BankMedBuffer}{}\renewcommand{\BankMedBuffer}{$-$0.220}
\providecommand{\BankMedPsi}{}\renewcommand{\BankMedPsi}{0.034}
\providecommand{\BankTrimSize}{}\renewcommand{\BankTrimSize}{1.282}
\providecommand{\BankExclSize}{}\renewcommand{\BankExclSize}{0.707}
\providecommand{\BankExclBuffer}{}\renewcommand{\BankExclBuffer}{$-$0.234}
\providecommand{\BankExclPsi}{}\renewcommand{\BankExclPsi}{0.023}
\providecommand{\BankExclIllSize}{}\renewcommand{\BankExclIllSize}{1.997}
\providecommand{\BankGapCutSize}{}\renewcommand{\BankGapCutSize}{67}
\providecommand{\BankGapCutBuffer}{}\renewcommand{\BankGapCutBuffer}{49}
\providecommand{\BankNlate}{}\renewcommand{\BankNlate}{105}
\providecommand{\BankNobsStag}{}\renewcommand{\BankNobsStag}{11{,}602}
\providecommand{\BankNobsDrop}{}\renewcommand{\BankNobsDrop}{8{,}575}
\providecommand{\BankPsiStag}{}\renewcommand{\BankPsiStag}{0.352}
\providecommand{\BankPsiDrop}{}\renewcommand{\BankPsiDrop}{0.367}
\providecommand{\BankMaxStagDiffSE}{}\renewcommand{\BankMaxStagDiffSE}{0.58}
\providecommand{\BankMaxDropDiffSE}{}\renewcommand{\BankMaxDropDiffSE}{1.14}
\providecommand{\BankImpRel}{}\renewcommand{\BankImpRel}{0.00}
\providecommand{\BankImpRelSE}{}\renewcommand{\BankImpRelSE}{0.00}
\providecommand{\BankSimStable}{}\renewcommand{\BankSimStable}{100.0}
\providecommand{\BankWidenInd}{}\renewcommand{\BankWidenInd}{1.25}
\providecommand{\BankWidenDir}{}\renewcommand{\BankWidenDir}{1.01}
\providecommand{\BankSizeIndDelta}{}\renewcommand{\BankSizeIndDelta}{[$-$0.004, 0.770]}
\providecommand{\BankSizeIndSim}{}\renewcommand{\BankSizeIndSim}{[0.071, 1.043]}


This section uses the data analysed by \citet{Kripfganz+Sarafidis:2025}: 350
US banks observed over 36 quarters (2006Q1--2014Q4), distributed with the
replication materials of their article. The dependent variable is the ratio
of non-performing loans to total loans (\code{NPL}), a measure of credit
risk. Operational inefficiency (\code{INEFF}) is treated as endogenous with
respect to the idiosyncratic error and instrumented by the ratio of interest
expenses to deposits (\code{INTEREST}); capital adequacy (\code{CAR}), size
(\code{SIZE}), the capital buffer (\code{BUFFER}), profitability
(\code{PROFIT}), asset quality (\code{QUALITY}) and liquidity
(\code{LIQUIDITY}) are exogenous with respect to the idiosyncratic error but
may be correlated with the factors. The row-standardized weights matrix
links each bank to the banks whose debt ratios are most strongly correlated
with its own. The weights file is an unlabelled matrix ordered like the
bank identifiers, so it is passed together with these identifiers:

\begin{CodeInput}
>>> banks = pd.read_stata("v113i06.dta")
>>> W = pd.read_csv("W.csv", header=None).to_numpy()
>>> ids = sorted(banks["ID"].unique())
>>> common = ["CAR", "SIZE", "BUFFER", "PROFIT", "QUALITY", "LIQUIDITY"]
>>> X, IV = ["INEFF"] + common, ["INTEREST"] + common
>>> spec = dict(W_ids=ids, splag=True, tlags=1, iv=IV, iv_lags=1,
...             iv_splags=True, std=True)
>>> res = sd.fit(banks, "NPL", X, "ID", "TIME", W, **spec)
>>> mg = sd.fit(banks, "NPL", X, "ID", "TIME", W, estimator="mg", **spec)
\end{CodeInput}

The specification has 28 instruments for 9 parameters. All results of this
section are produced by \code{replication/bank\_replication.py}.

\subsection{Cross-software check}

Table~\ref{tab:bank_compare} compares the estimates with those printed in
Tables~1 and~5 of \citet{Kripfganz+Sarafidis:2025}. For both the 2SIV and
the mean-group estimator, the coefficients and standard errors agree with
the published ones in every printed digit: the largest absolute difference
is below $10^{-6}$, the precision of the printed numbers. The $J$ statistic
(\BankJ), the factor and idiosyncratic standard deviations, the numbers of
factors, and the long-run direct, indirect and total impacts of Table~3 of
the original article are reproduced to the same precision. A test in the
package (\code{tests/test\_stata\_replication.py}) checks this agreement
whenever the replication materials are available.

\begin{table}[t!]
\centering\scriptsize
\begin{tabular}{lrrrrrrrr}
\toprule
 & \multicolumn{4}{c}{2SIV (Table~1)} & \multicolumn{4}{c}{Mean group (Table~5)} \\
\cmidrule(lr){2-5}\cmidrule(lr){6-9}
 & \multicolumn{2}{c}{Coefficient} & \multicolumn{2}{c}{Std.\ error} & \multicolumn{2}{c}{Coefficient} & \multicolumn{2}{c}{Std.\ error} \\
\cmidrule(lr){2-3}\cmidrule(lr){4-5}\cmidrule(lr){6-7}\cmidrule(lr){8-9}
 & \pkg{spdfiv} & \proglang{Stata} & \pkg{spdfiv} & \proglang{Stata} & \pkg{spdfiv} & \proglang{Stata} & \pkg{spdfiv} & \proglang{Stata} \\
\midrule
\code{W.NPL} & 0.3943 & 0.3943 & 0.0849 & 0.0849 & 0.0316 & 0.0316 & 0.0511 & 0.0511 \\
\code{L1.NPL} & 0.2899 & 0.2899 & 0.0544 & 0.0544 & 0.3005 & 0.3005 & 0.0149 & 0.0149 \\
\code{INEFF} & 0.4474 & 0.4474 & 0.1046 & 0.1046 & 0.7588 & 0.7588 & 0.1584 & 0.1584 \\
\code{CAR} & 0.0305 & 0.0305 & 0.0058 & 0.0058 & 0.2181 & 0.2181 & 0.0263 & 0.0263 \\
\code{SIZE} & 0.2226 & 0.2226 & 0.0942 & 0.0942 & 2.0040 & 2.0040 & 0.3385 & 0.3385 \\
\code{BUFFER} & $-$0.0545 & $-$0.0545 & 0.0119 & 0.0119 & $-$0.3764 & $-$0.3764 & 0.0420 & 0.0420 \\
\code{PROFIT} & $-$0.0053 & $-$0.0053 & 0.0018 & 0.0018 & $-$0.0180 & $-$0.0180 & 0.0059 & 0.0059 \\
\code{QUALITY} & 0.1830 & 0.1830 & 0.0308 & 0.0308 & 0.2873 & 0.2873 & 0.1387 & 0.1387 \\
\code{LIQUIDITY} & 2.4524 & 2.4524 & 0.2696 & 0.2696 & 6.3302 & 6.3302 & 0.5059 & 0.5059 \\
\midrule
$J$ statistic & 18.825 & 18.825 & & & & & & \\
$\sigma_f$ & 0.6416 & 0.6416 & & & & & & \\
$\sigma_e$ & 0.9038 & 0.9038 & & & & & & \\
\bottomrule
\end{tabular}

\caption{Bank data: estimates from \pkg{spdfiv} and from \code{spxtivdfreg}
as published in Tables~1 and~5 of \citet{Kripfganz+Sarafidis:2025}.}
\label{tab:bank_compare}
\end{table}

Reaching this agreement required five implementation details of
\code{spxtivdfreg} that are not stated in the equations of the article and
that we identified from the source code of the command, distributed under
the GPL with the article. (i) With option \code{std}, variables are
standardized across banks \emph{within each period} before the factors are
extracted. (ii) The eigenvalue-ratio criterion chooses between one and
\code{factmax} factors, without the mock eigenvalue of
\citet{Ahn+Horenstein:2013}. (iii) All lag blocks of a set of instruments
use the largest number of factors selected for any of them. (iv) The
variance of the 2SIV estimator and the $J$ statistic use the weight matrix
computed from the first-stage residuals, as in Equation~(11) of the article.
(v) The instruments are demeaned before, but not after, the factors are
projected out. Point (v) leaves $\hat{\bm Z}_i^\top\bm C_i$ and
$\hat{\bm Z}_i^\top\bm y_i$ unchanged, because the regressors are demeaned,
but it changes the IV weight $\hat{\bm Z}_i^\top\hat{\bm Z}_i$; its effect on
the pooled estimator is of the order $10^{-4}$, whereas for the mean-group
estimator, which uses a separate weight for each bank, it changes the
spatial coefficient from 0.087 to 0.032. Points (i)--(v) are the defaults
of \pkg{spdfiv}; points (i), (ii) and (iv) can be changed through arguments
of \fct{fit} (Table~\ref{tab:options}). Without them, the selected numbers
of factors differ (one instead of two in the instruments, none instead of
one in the residuals), and the coefficient on \code{LIQUIDITY} is 1.50
instead of 2.45. Making such details explicit is, in our view, a useful
by-product of an independent implementation: they matter for the results
but cannot be inferred from the published equations.

\subsection{Long-run impacts}

Table~\ref{tab:bank_impacts} reports the long-run impacts with Delta-method
and simulation intervals (10{,}000 draws, all of which satisfy the stability
condition). The two methods agree for the direct impacts: the simulation
intervals are on average as wide as the Delta intervals, with a median width
ratio of \BankWidenDir. They disagree for the indirect impacts, which
depend on the multiplier $1/(1-\rho_1-\psi_0)$ with $\hat\rho_1+\hat\psi_0=0.68$.
The simulation intervals are skewed to the right, and are \BankWidenInd{}
times as wide as the Delta intervals in the median. The disagreement changes
a conclusion: the Delta interval for the indirect impact of \code{SIZE},
\BankSizeIndDelta, contains zero, whereas the simulation interval,
\BankSizeIndSim, does not. Given the Monte Carlo evidence of
Section~\ref{sec:mc}, we read this as showing that the uncertainty about
the spillovers is asymmetric, not as establishing that the simulation
interval is the more accurate one in this sample.

\begin{table}[t!]
\centering\scriptsize
\begin{tabular}{llrrrrr}
\toprule
 & & & \multicolumn{2}{c}{Delta method} & \multicolumn{2}{c}{Simulation} \\
\cmidrule(lr){4-5}\cmidrule(lr){6-7}
Impact & Variable & Estimate & s.e. & 95\% interval & s.e. & 95\% interval \\
\midrule
direct & \code{INEFF} & 0.647 & 0.159 & [0.335, 0.959] & 0.163 & [0.339, 0.979] \\
 & \code{CAR} & 0.044 & 0.009 & [0.026, 0.062] & 0.009 & [0.027, 0.063] \\
 & \code{SIZE} & 0.322 & 0.142 & [0.044, 0.600] & 0.143 & [0.056, 0.614] \\
 & \code{BUFFER} & $-$0.079 & 0.018 & [$-$0.115, $-$0.043] & 0.018 & [$-$0.117, $-$0.044] \\
 & \code{PROFIT} & $-$0.008 & 0.002 & [$-$0.012, $-$0.003] & 0.002 & [$-$0.012, $-$0.003] \\
 & \code{QUALITY} & 0.265 & 0.047 & [0.173, 0.356] & 0.047 & [0.179, 0.363] \\
 & \code{LIQUIDITY} & 3.547 & 0.445 & [2.674, 4.420] & 0.462 & [2.736, 4.513] \\
\addlinespace
indirect & \code{INEFF} & 0.769 & 0.335 & [0.112, 1.427] & 0.588 & [0.294, 1.909] \\
 & \code{CAR} & 0.052 & 0.024 & [0.006, 0.099] & 0.042 & [0.020, 0.138] \\
 & \code{SIZE} & 0.383 & 0.198 & [$-$0.004, 0.770] & 0.285 & [0.071, 1.043] \\
 & \code{BUFFER} & $-$0.094 & 0.043 & [$-$0.178, $-$0.010] & 0.074 & [$-$0.244, $-$0.035] \\
 & \code{PROFIT} & $-$0.009 & 0.005 & [$-$0.018, 0.000] & 0.007 & [$-$0.024, $-$0.002] \\
 & \code{QUALITY} & 0.315 & 0.141 & [0.039, 0.591] & 0.244 & [0.121, 0.820] \\
 & \code{LIQUIDITY} & 4.218 & 1.742 & [0.803, 7.633] & 3.100 & [1.801, 10.548] \\
\addlinespace
total & \code{INEFF} & 1.417 & 0.427 & [0.579, 2.254] & 0.659 & [0.724, 2.731] \\
 & \code{CAR} & 0.097 & 0.029 & [0.039, 0.154] & 0.046 & [0.052, 0.191] \\
 & \code{SIZE} & 0.705 & 0.310 & [0.097, 1.312] & 0.380 & [0.140, 1.529] \\
 & \code{BUFFER} & $-$0.173 & 0.054 & [$-$0.279, $-$0.066] & 0.083 & [$-$0.343, $-$0.089] \\
 & \code{PROFIT} & $-$0.017 & 0.006 & [$-$0.029, $-$0.004] & 0.008 & [$-$0.034, $-$0.006] \\
 & \code{QUALITY} & 0.580 & 0.167 & [0.252, 0.907] & 0.265 & [0.334, 1.136] \\
 & \code{LIQUIDITY} & 7.765 & 1.904 & [4.034, 11.496] & 3.262 & [5.021, 14.366] \\
\bottomrule
\end{tabular}

\caption{Bank data: long-run average impacts with 95\% intervals from the
Delta method and from Krinsky--Robb simulation (10{,}000 draws).}
\label{tab:bank_impacts}
\end{table}

\subsection{What drives the mean-group estimates?}

\citet{Kripfganz+Sarafidis:2025} found that the MG estimates of some slopes
differ from the 2SIV estimates by an order of magnitude, although the $J$
test does not reject the homogeneous model, and that the spatial spillover
disappears under heterogeneous slopes; they suggested that a few outlying
banks might be responsible. The diagnostics of Section~\ref{sec:diag}
allow this conjecture to be examined.

\begin{table}[t!]
\centering\small
\begin{tabular}{rlrrrr}
\toprule
Bank & Coefficient & $\hat\theta_i$ & Modified $z$ & Change in MG (s.e.) & cond$(\bm Z_i^\top\bm Z_i)$ \\
\midrule
147 & \code{BUFFER} & 3.618 & 7.58 & 0.272 & $1.7\times 10^{6}$ \\
26 & \code{SIZE} & 32.324 & 9.14 & 0.257 & $1.7\times 10^{6}$ \\
112 & \code{SIZE} & 30.991 & 8.76 & 0.245 & $3.7\times 10^{6}$ \\
296 & \code{SIZE} & 30.810 & 8.70 & 0.244 & $2.5\times 10^{7}$ \\
276 & \code{BUFFER} & 2.726 & 5.82 & 0.212 & $1.0\times 10^{7}$ \\
199 & \code{BUFFER} & $-$3.368 & $-$6.22 & $-$0.204 & $1.9\times 10^{7}$ \\
79 & \code{SIZE} & 25.680 & 7.22 & 0.200 & $8.6\times 10^{5}$ \\
16 & \code{BUFFER} & $-$3.207 & $-$5.90 & $-$0.193 & $1.2\times 10^{7}$ \\
\bottomrule
\end{tabular}

\caption{Bank data: the eight banks with the largest influence on the MG
estimates of the \code{SIZE} and \code{BUFFER} slopes. ``Change in MG'' is
the change in the MG estimate from deleting the bank, in units of its
standard error.}
\label{tab:bank_influence}
\end{table}

\begin{table}[t!]
\centering\scriptsize
\begin{tabular}{lcccc}
\toprule
 & MG & Trimmed (10\%) & Median & MG excl.\ flagged \\
\midrule
\code{W.NPL} & 0.032 & 0.020 & 0.034 & 0.023 \\
\footnotesize  & \footnotesize [$-$0.069, 0.132] & \footnotesize [$-$0.063, 0.102] & \footnotesize [$-$0.031, 0.117] & \footnotesize [$-$0.065, 0.109] \\
\code{L1.NPL} & 0.301 & 0.301 & 0.301 & 0.303 \\
\footnotesize  & \footnotesize [0.271, 0.330] & \footnotesize [0.272, 0.334] & \footnotesize [0.268, 0.345] & \footnotesize [0.271, 0.334] \\
\code{INEFF} & 0.759 & 0.507 & 0.313 & 0.191 \\
\footnotesize  & \footnotesize [0.448, 1.069] & \footnotesize [0.271, 0.763] & \footnotesize [0.099, 0.572] & \footnotesize [$-$0.001, 0.403] \\
\code{CAR} & 0.218 & 0.184 & 0.123 & 0.133 \\
\footnotesize  & \footnotesize [0.167, 0.270] & \footnotesize [0.143, 0.225] & \footnotesize [0.094, 0.156] & \footnotesize [0.094, 0.171] \\
\code{SIZE} & 2.004 & 1.282 & 0.807 & 0.707 \\
\footnotesize  & \footnotesize [1.341, 2.668] & \footnotesize [0.796, 1.750] & \footnotesize [0.232, 1.198] & \footnotesize [0.300, 1.123] \\
\code{BUFFER} & $-$0.376 & $-$0.321 & $-$0.220 & $-$0.234 \\
\footnotesize  & \footnotesize [$-$0.459, $-$0.294] & \footnotesize [$-$0.389, $-$0.254] & \footnotesize [$-$0.281, $-$0.167] & \footnotesize [$-$0.301, $-$0.171] \\
\code{PROFIT} & $-$0.018 & $-$0.016 & $-$0.008 & $-$0.020 \\
\footnotesize  & \footnotesize [$-$0.030, $-$0.006] & \footnotesize [$-$0.025, $-$0.007] & \footnotesize [$-$0.018, $-$0.001] & \footnotesize [$-$0.030, $-$0.011] \\
\code{QUALITY} & 0.287 & 0.179 & 0.081 & 0.159 \\
\footnotesize  & \footnotesize [0.015, 0.559] & \footnotesize [0.021, 0.331] & \footnotesize [$-$0.055, 0.214] & \footnotesize [0.011, 0.304] \\
\code{LIQUIDITY} & 6.330 & 4.921 & 3.200 & 3.383 \\
\footnotesize  & \footnotesize [5.339, 7.322] & \footnotesize [4.075, 5.772] & \footnotesize [2.573, 4.049] & \footnotesize [2.691, 4.033] \\
\midrule
Units & 350 & 350 & 350 & 255 \\
\bottomrule
\end{tabular}

\caption{Bank data: MG estimates and robust summaries of the bank-specific
estimates, with 95\% intervals (normal for MG, bootstrap percentile with 999
resamples of banks for the others).}
\label{tab:bank_mg}
\end{table}

Table~\ref{tab:bank_influence} shows that single banks move the MG estimate
of \code{SIZE} or \code{BUFFER} by about 0.19 to 0.27 standard errors. Their
instrument matrices are far from singular (condition numbers below
$10^{8}$), so these are not numerical artefacts. The flags mark
\BankNflagSize{} banks for \code{SIZE}, \BankNflagBuffer{} for
\code{BUFFER}, and \BankNflag{} of the \BankNunits{} banks for at least one
covariate slope. The distribution of the bank-specific estimates is
therefore heavy tailed rather than contaminated by a handful of outliers.

Table~\ref{tab:bank_mg} gives the robust summaries. For \code{SIZE} the
median, \BankMedSize, closes \BankGapCutSize\% of the gap between the MG
estimate (\BankMGsize) and the 2SIV estimate (\BankTwoSize); for
\code{BUFFER} the median closes \BankGapCutBuffer\% of the gap
(\BankMedBuffer{} against \BankMGbuffer{} and \BankTwoBuffer). Excluding the
flagged banks gives similar values (\BankExclSize{} and \BankExclBuffer).
The spatial coefficient behaves differently: the median (\BankMedPsi) and
the MG estimate without the flagged banks (\BankExclPsi) are as small as the
MG estimate itself (\BankMGpsi) and far below the 2SIV estimate
(\BankTwoPsi). Five banks have nearly singular instrument matrices, because
\code{QUALITY} does not vary over time for them, but removing them leaves
the MG estimates practically unchanged (\code{SIZE}: \BankExclIllSize).

The conjecture is thus partly supported. Extreme bank-level estimates
account for about two thirds of the difference in the \code{SIZE} slope and
for about half of it in the \code{BUFFER} slope, but not for the
disappearance of the spatial spillover, which is a feature of the bulk of
the distribution. The MG estimate of the spatial coefficient is also
sensitive to the construction of the instruments: without double
defactoring it is \BankMGpsiNoDD{} instead of \BankMGpsi{}, and point (v)
above moves it as much. One explanation consistent with this evidence,
which we do not pursue here, is the finite-sample behaviour of bank-level
IV estimation with 28 instruments and 35 observations per bank, a ratio at
which IV estimates are known to be pulled towards least squares.

\subsection{Staggered entry}

Finally, the bank panel is made unbalanced artificially, so that the answer
is known: for a random \BankNlate{} of the 350 banks the first 1 to 12
quarters are deleted. Table~\ref{tab:bank_staggered} compares three
estimates: the full balanced panel, the experimental unbalanced estimator
using all \BankNobsStag{} remaining observations, and the balanced
alternative that drops the late banks (\BankNobsDrop{} observations).

\begin{table}[t!]
\centering\scriptsize
\begin{tabular}{lrrrrrr}
\toprule
 & \multicolumn{2}{c}{Full panel} & \multicolumn{2}{c}{Staggered, all banks} & \multicolumn{2}{c}{Late banks dropped} \\
\cmidrule(lr){2-3}\cmidrule(lr){4-5}\cmidrule(lr){6-7}
 & Coef. & s.e. & Coef. & s.e. & Coef. & s.e. \\
\midrule
\code{W.NPL} & 0.394 & 0.085 & 0.352 & 0.106 & 0.367 & 0.105 \\
\code{L1.NPL} & 0.290 & 0.054 & 0.305 & 0.054 & 0.243 & 0.058 \\
\code{INEFF} & 0.447 & 0.105 & 0.435 & 0.101 & 0.456 & 0.117 \\
\code{CAR} & 0.031 & 0.006 & 0.032 & 0.006 & 0.030 & 0.006 \\
\code{SIZE} & 0.223 & 0.094 & 0.199 & 0.105 & 0.330 & 0.124 \\
\code{BUFFER} & $-$0.055 & 0.012 & $-$0.051 & 0.013 & $-$0.055 & 0.013 \\
\code{PROFIT} & $-$0.005 & 0.002 & $-$0.004 & 0.002 & $-$0.006 & 0.002 \\
\code{QUALITY} & 0.183 & 0.031 & 0.173 & 0.031 & 0.182 & 0.029 \\
\code{LIQUIDITY} & 2.452 & 0.270 & 2.525 & 0.282 & 2.468 & 0.325 \\
\midrule
Observations & \multicolumn{2}{c}{12{,}250} & \multicolumn{2}{c}{11{,}602} & \multicolumn{2}{c}{8{,}575} \\
Banks & \multicolumn{2}{c}{350} & \multicolumn{2}{c}{350} & \multicolumn{2}{c}{245} \\
\bottomrule
\end{tabular}

\caption{Bank data with artificial staggered entry: 2SIV estimates.}
\label{tab:bank_staggered}
\end{table}

The unbalanced estimates stay within \BankMaxStagDiffSE{} standard errors
of the full-panel estimates. Dropping the late banks moves the estimates by
up to \BankMaxDropDiffSE{} standard errors, because it discards 30\% of the
network as well as the observations. In this application the experimental
option thus recovers more of the full-sample information than the balanced
alternative available in existing software.
